\subsection{Uniform dyadic support matching in every weight}

\begin{result}{E.158}{Proved}{The dyadic localization map on the ordinary support pattern}
Put
\[
 \Sigma=\{2,23,41\}\cup\mathcal T,\qquad
 \mathcal T\subset\{q\text{ prime}:q\equiv5\pmod8\}.
 \tag{E.158.1}
\]
Fix compatible Tate bases. For each odd \(w\ge3\), there is a
generator \(\epsilon_w\) of
\(\ker[H^1(G_{\mathbb Q},\mathbb Z_2(w))
\to H^1(G_{\mathbb R},\mathbb Z_2(w))]\)
whose reduction modulo eight, under
\(\mathbb Z/8(w)=\mu_8\), is \(\kappa_8(2)\).
Let \(\eta_w=(\chi_2^w-1)/2\). Then
\[
 H^1(G_{\mathbb Q},\mathbb Z_2(w))
       =\mathbb Z_2\epsilon_w\oplus\mathbb F_2\eta_w.
 \tag{E.158.2}
\]
Write \(e_{q,w}=\operatorname{loc}_q(\epsilon_w)\). The
localization table is
\[
\begin{array}{c|c|c|c}
 \text{place}&H^1(\mathbb Q_q,\mathbb Z_2(w))
     &\operatorname{loc}(\epsilon_w)&\operatorname{loc}(\eta_w)\\\hline
 \mathbb R&\mathbb F_2&0&1\\
 23&\mathbb Z/2&0&1\\
 41&\mathbb Z/8&e_{41,w}\text{ of order }4&2e_{41,w}\\
 q\in\mathcal T&\mathbb Z/4&e_{q,w}\text{ a generator}&2e_{q,w}.
\end{array}
\tag{E.158.3}
\]
The entries in the last two rows are determined by the one
finite Kummer class \(\kappa_8(2)\), independently of \(w\).
In particular the joint image on the real place and all odd
support places is \(\mathbb F_2\oplus\mathbb Z/4\), and its
kernel is \(4\mathbb Z_2\epsilon_w\). Explicitly,
\(v\epsilon_w+\tau\eta_w\) maps to \(\tau\) at the real place
and at twenty-three, and to
\[
                     (v+2\tau)e_{q,w}
              \quad(q=41\text{ or }q\in\mathcal T).
 \tag{E.158.4}
\]

For even \(w\ge2\), put \(a_w=2+v_2(w)\) and
\(\eta_w=(\chi_2^w-1)/2^{a_w}\). The global group is
\(\mathbb Z/2^{a_w}\eta_w\). Its localization is injective
at every odd prime, is an isomorphism at each
\(q\in\mathcal T\), and has index two at twenty-three
and forty-one. Its real localization is zero.
\end{result}
\begin{proof}
E.155--E.156 give the global groups, their real restrictions,
and their equality with the groups unramified outside
\(\{2,\infty\}\). For odd \(w\), every unit modulo eight
has square one, so \(\chi_2^w=\chi_2\) modulo eight.
The coefficient sequence and E.155 give the surjection
\[
 H^1(G_{\mathbb Q,\{2,\infty\}},\mathbb Z_2(w))
       \longrightarrow
 H^1(G_{\mathbb Q,\{2,\infty\}},\mu_8).
 \tag{E.158.5}
\]
The Kummer class of two belongs to its target. Its real
restriction is zero, and
\(H^1(\mathbb R,\mathbb Z_2(w))\to H^1(\mathbb R,\mu_8)\)
is an isomorphism. A lift therefore has real restriction
zero. Its reduction modulo two is \(\kappa_2(2)\), so
E.156 shows that it generates the free real kernel.
This constructs \(\epsilon_w\) with the stated normalization;
there is no assumption that an arbitrary mod-eight class
lifts to every Tate weight.

For an odd prime \(q\) and \(w\ge2\), the integral inertia
part of a Tate cocycle is zero, since it is annihilated by
the nonzero integer \(q^{w-1}-1\). Thus
\[
 H^1(\mathbb Q_q,\mathbb Z_2(w))
       =\mathbb Z_2/(q^w-1).
 \tag{E.158.6}
\]
For odd \(w\), \(v_2(q^w-1)=v_2(q-1)\). On the primes
in E.158.1 this valuation is at most three, so reduction
modulo eight is injective on the integral local group.
Its \(\epsilon_w\)-class is consequently determined
exactly by the unramified local class \(\kappa_8(2)\).

At twenty-three, \(2^{11}\equiv1\pmod{23}\), so this
class is zero in the local group of order two.
At forty-one, \(2^{10}\equiv-1\pmod{41}\) and
\(2^4\not\equiv1\pmod{41}\); hence two has order twenty
in \(\mathbb F_{41}^{\times}\). Its class modulo eighth
powers has order four. Equivalently,
\(2^5=32\pmod{41}\) has square \(-1\).
At \(q=8k+5\), Gauss's lemma shows that two is a
nonsquare: among the least residues of \(2,4,\ldots,
2(4k+2)\), exactly \(2k+1\) exceed \(q/2\).
Its class modulo fourth powers therefore has order four.
These calculations prove the \(\epsilon_w\)-column.

The real value of \(\eta_w\) is \(-1\), representing
the nonzero class of \(\mathbb F_2\). Its Frobenius
value is \((q^w-1)/2\). At twenty-three this is odd;
at forty-one it is four modulo eight, since
\(41^w\equiv9\pmod{16}\); at each \(q\in\mathcal T\)
it is two modulo four. These are respectively the
claimed classes \(1,2e_{41,w},2e_{q,w}\).
The real place distinguishes \(\tau\), and the
forty-one row distinguishes \(v+2\tau\) modulo four.
This proves the exact image and kernel, including when
\(\mathcal T\) is empty.

For even \(w\), the elementary lifting-the-exponent
formula gives
\[
 v_2(q^w-1)=v_2(q-1)+v_2(q+1)+v_2(w)-1.
 \tag{E.158.7}
\]
It is at least \(a_w\) for every odd prime. The local
value \((q^w-1)/2^{a_w}\) has exact order \(2^{a_w}\)
in E.158.6, proving injectivity. At \(q\equiv5\pmod8\)
the valuation is \(a_w\); at twenty-three and forty-one
it is \(a_w+1\). These give the stated indices.
Finally the even-weight real first group is zero.
\end{proof}

\begin{result}{E.159}{Proved}{An exact finite parameter test for every dyadic matching step}
Let a fixed ordinary completed pair have support E.158.1.
Write \(N_j=\Pi_2/\Gamma_{j+1}\Pi_2\) and
\(L_w=\Gamma_w\Pi_2/\Gamma_{w+1}\Pi_2
\simeq\mathbb Z_2(w)^{d_w}\).
Suppose an integral global \(N_{w-1}\)-section is matched
to its ordinary boundary point at every odd support
place and to its ordinary real class.
It has an integral global \(N_w\)-lift \(s_w\).

Fix such a lift. At each odd support place let
\(\delta_q\in H^1(\mathbb Q_q,L_w)\) be the correction
needed to match the ordinary point lift. These classes
are independent of the lower aligning frame. At the
real place choose one aligning frame and let
\(\xi_\infty\in H^1(\mathbb R,L_w)\) be its correction.
For odd \(w\), define the finite affine subset
\[
 \mathcal R_\infty=\xi_\infty+
 \operatorname{im}\left[
 H^0(\mathbb R,N_{w-1}^{s_{w-1}})
       \longrightarrow H^1(\mathbb R,L_w)\right]
 \subset\mathbb F_2^{d_w}.
 \tag{E.159.1}
\]
The boundary map uses the actual \(N_w\)-extension
twisted by \(s_w\). It records all real aligning frames.
For matching relative to the chosen real frame alone,
replace \(\mathcal R_\infty\) by the singleton
\(\{\xi_\infty\}\).

For odd \(w\), simultaneous matching at all odd support
places and at the real place is equivalent to the
following three finite-parameter conditions:
\[
\begin{gathered}
 \delta_{23}\in\mathcal R_\infty,\qquad
 \delta_{41}\in
       (\mathbb Z/4)e_{41,w}^{\,d_w},\\
 h=e_{41,w}^{-1}\delta_{41}\in(\mathbb Z/4)^{d_w},
 \qquad
 e_{q,w}^{-1}\delta_q=h
                       \quad(q\in\mathcal T).
\end{gathered}
\tag{E.159.2}
\]
Here multiplication and the inverse notation are coordinatewise,
and the inverse in the forty-one row is defined only on the
displayed order-four subgroup. If these conditions hold, all
matching global corrections are
\[
 v\epsilon_w+\tau\eta_w,\qquad
 \tau=\delta_{23},\qquad
 v\equiv h-2\tau\pmod4.
 \tag{E.159.3}
\]
Thus the matching lift fibre, when nonempty, is a torsor
under \(4\mathbb Z_2^{d_w}\).

For even \(w\), simultaneous matching is equivalent to
\(\delta_{41}\) belonging to the image of the global
\((\mathbb Z/2^{a_w})^{d_w}\)-torsor and to its unique
preimage localizing to \(\delta_q\) at every other odd
support place. The matching lift, if it exists, is unique.
The real match imposes no further first-cohomology test
in that parity.

In either parity, inconsistency has a witness using at
most two odd support places, together with the real
place when required. These are exact tests for a fixed
integral predecessor. They do not assert that a passing
predecessor or an infinite passing branch always exists,
and they impose no full point-class matching at the
coefficient prime two.
\end{result}
\begin{proof}
For \(w=2\), A.25 supplies the global class-two lift.
For \(w\ge3\), B.25 puts the next global obstruction in the
fixed-ramification dyadic second group; E.155 identifies
it with its real restriction. The chosen real point
class supplies a local lift, so the obstruction is zero.
This proves the existence of \(s_w\), before any assertion
about retaining all matches.

The global lifts above the fixed predecessor form the
usual \(H^1(G_{\mathbb Q},L_w)\)-torsor. There is no
global lower stabilizer: a nontrivial invariant element
has a first nonzero weight \(i\), on which the action is
\(\chi_2^i\); an element with cyclotomic value five would
annihilate that nonzero torsion-free coordinate by
\(5^i-1\). The same argument at an odd prime uses
arithmetic Frobenius and \(q^i-1\), and proves that its
lower stabilizer is trivial. Twisting by the section
does not change this graded action. Consequently the
local corrections \(\delta_q\) are well-defined and the
local central torsor actions are free.

At the real place that stabilizer argument is not
available: even weights have invariant coordinates.
Changing the chosen lower frame by an invariant lower
element changes the central comparison by precisely the
boundary class in E.159.1. Its image is a subgroup of
the finite first-cohomology group; the connecting map
is additive because the kernel \(L_w\) is central.
Thus E.159.1 is exactly the set of real central
corrections giving the unframed ordinary point class.
This step retains, rather than discards, the real
aligning-frame ambiguity.

Apply E.158 in each of the \(d_w\) central coordinates.
For odd \(w\), the twenty-three row forces
\(\tau=\delta_{23}\). Real matching is therefore exactly
the first condition of E.159.2. The forty-one row says
that the correction must lie in its order-four subgroup
and forces \(v+2\tau=h\) modulo four. Every remaining
support row now has the same scalar \(h\), yielding the
last conditions of E.159.2. These calculations prove
necessity, sufficiency, and the entire solution set
E.159.3. Its translation kernel is \(4\mathbb Z_2^{d_w}\).
For even \(w\), E.158 makes the forty-one localization
injective on the entire finite global torsor, so it
either fails alone or determines the unique correction.
The local real first group is zero, so a lower real
point match automatically persists once a lift exists.

For the witness assertion, an odd failure of the real
condition uses twenty-three and the real place; a
failure of subgroup membership uses forty-one alone.
Otherwise forty-one determines \(h\), and any failed
remaining equation is witnessed by forty-one and that
one support place. The even argument is identical
after its injective forty-one test. Infinite support
introduces no additional compactness step in this
criterion: a fixed candidate scalar either satisfies
each displayed local equation or fails a particular
one. Existence of candidates at all subsequent depths
remains a separate problem.
\end{proof}
